\section{Sensitivity Analysis}
\label{sec:sensitivity}

Comprehensive sensitivity analysis is essential for validating model robustness and identifying critical parameters. We conduct multi-factor sensitivity analysis examining parameter sensitivity, assumption sensitivity, and algorithm sensitivity.

\subsection{Parameter Sensitivity Analysis}

\subsubsection{Reliability Threshold $\alpha$}

We vary the reliability threshold $\alpha \in \{0.70, 0.75, 0.80, 0.85, 0.90, 0.95\}$:

\begin{table}[H]
\centering
\caption{Sensitivity to Reliability Threshold $\alpha$}
\begin{tabular}{lcccccc}
\toprule
$\alpha$ & Expected Time (hrs) & Std Dev & Actual Reliability & Feasible & Time Change & Cost Change \\
\midrule
0.70 & 32.8 & 9.4 & 0.74 & Yes & -10.1\% & -8.3\% \\
0.75 & 34.1 & 8.7 & 0.79 & Yes & -6.6\% & -5.1\% \\
0.80 & 35.4 & 8.1 & 0.84 & Yes & -3.0\% & -2.4\% \\
0.85 & 36.5 & 7.6 & 0.89 & Yes & \textbf{Baseline} & \textbf{Baseline} \\
0.90 & 38.9 & 6.9 & 0.92 & Yes & +6.6\% & +7.8\% \\
0.95 & 43.7 & 5.8 & 0.96 & Yes & +19.7\% & +18.2\% \\
\bottomrule
\end{tabular}
\end{table}

\textbf{Key Insights:}
\begin{itemize}
\item Marginal cost of increasing $\alpha$ from 0.85 to 0.90: +6.6\% time, +7.8\% cost
\item High-reliability solutions ($\alpha \geq 0.90$) become increasingly expensive
\item $\alpha = 0.85$ represents reasonable balance between reliability and efficiency
\item Actual reliability consistently exceeds threshold (conservative bias in model)
\end{itemize}

\subsubsection{Fleet Size (Number of Vehicles $K$)}

We vary fleet size $K \in \{4, 6, 8, 10, 12\}$ for instance I-Medium-1 ($N=30$):

\begin{table}[H]
\centering
\caption{Sensitivity to Fleet Size $K$}
\begin{tabular}{lccccc}
\toprule
$K$ (Vehicles) & Time (hrs) & Utilization & Reliability & Cost per Delivery & Diminishing Returns \\
\midrule
4 & 42.3 & 96\% & 0.84 & \$89.2 & Overloaded \\
6 & 36.5 & 89\% & 0.89 & \$78.4 & \textbf{Optimal} \\
8 & 34.1 & 72\% & 0.91 & \$92.3 & Underutilized \\
10 & 33.2 & 61\% & 0.92 & \$108.7 & Poor efficiency \\
12 & 32.8 & 54\% & 0.92 & \$124.1 & Poor efficiency \\
\bottomrule
\end{tabular}
\end{table}

\textbf{Key Insights:}
\begin{itemize}
\item Minimum fleet size for feasibility: $K = 5$ vehicles
\item Optimal fleet size: $K = 6$ (balances utilization and reliability)
\item Diminishing returns beyond $K = 8$ (utilization $< 70\%$)
\item Humanitarian organizations should maintain 20\% spare capacity for emergencies
\end{itemize}

\subsubsection{Vehicle Capacity $C$}

We vary vehicle capacity $C \in \{50, 75, 100, 125, 150\}$ units:

\begin{table}[H]
\centering
\caption{Sensitivity to Vehicle Capacity $C$}
\begin{tabular}{lcccc}
\toprule
Capacity $C$ & Routes Required & Time (hrs) & Reliability & Utilization \\
\midrule
50 & 12 & 48.7 & 0.82 & 94\% \\
75 & 8 & 39.2 & 0.87 & 88\% \\
100 & 6 & 36.5 & 0.89 & 82\% \\
125 & 5 & 35.1 & 0.88 & 79\% \\
150 & 4 & 34.8 & 0.86 & 76\% \\
\bottomrule
\end{tabular}
\end{table}

\textbf{Key Insights:}
\begin{itemize}
\item Capacity beyond $C = 100$ provides minimal benefit (diminishing returns)
\item Larger vehicles reduce number of routes but increase individual route risk
\item Recommended capacity: $C = 100$ units (standard humanitarian vehicle configuration)
\end{itemize}

\subsection{Assumption Sensitivity Analysis}

We test robustness to violations of modeling assumptions made in Section \ref{sec:problem_analysis}.

\subsubsection{Travel Time Independence}

Our model assumes travel times on different segments are independent. We test correlated travel times:

\begin{table}[H]
\centering
\caption{Sensitivity to Travel Time Correlation}
\begin{tabular}{lcccc}
\toprule
Correlation $\rho$ & Expected Time & Std Dev & Reliability & Model Error \\
\midrule
0.00 (Independent) & 36.5 & 7.6 & 0.89 & 0.0\% \\
0.20 (Weak) & 36.8 & 8.1 & 0.88 & +0.8\% \\
0.40 (Moderate) & 37.4 & 8.9 & 0.86 & +2.5\% \\
0.60 (Strong) & 38.7 & 10.2 & 0.83 & +6.0\% \\
\bottomrule
\end{tabular}
\end{table}

\textbf{Key Insights:}
\begin{itemize}
\item Model robust to weak-moderate correlation (error $< 3\%$)
\item Strong correlation ($\rho = 0.60$) reduces reliability by 6 percentage points
\item Somalia data shows $\rho \approx 0.34$ (moderate correlation)
\item Model underestimates variance by 2.5\%—acceptable for planning purposes
\end{itemize}

\subsubsection{Stationary Travel Time Distributions}

Our model assumes travel time distributions do not change within planning horizon. We test non-stationary distributions (simulating sudden weather change):

\begin{table}[H]
\centering
\caption{Sensitivity to Non-Stationary Distributions}
\begin{tabular}{lcccc}
\toprule
Scenario & Mean Change & Std Dev Change & Reliability & Adaptive Needed? \\
\midrule
Stationary (baseline) & 0\% & 0\% & 0.89 & No \\
Gradual change (10\%/hr) & +15\% & +20\% & 0.84 & Maybe \\
Sudden change (at $t=4$ hrs) & +30\% & +40\% & 0.76 & \textbf{Yes} \\
\bottomrule
\end{tabular}
\end{tabular}

\textbf{Key Insights:}
\begin{itemize}
\item Model handles gradual changes (seasonal transitions) well
\item Sudden changes (flash floods) significantly impact reliability
\item Recommendation: Implement dynamic re-optimization for extreme weather events
\item Future work: Integrate weather forecasting for proactive adaptation
\end{itemize}

\subsubsection{Demand Variability}

Our model assumes known deterministic demands. We test demand uncertainty ($\pm 20\%$ variation):

\begin{table}[H]
\centering
\caption{Sensitivity to Demand Uncertainty}
\begin{tabular}{lcccc}
\toprule
Demand Variation & Time & Feasibility & Cost Impact & Recommendation \\
\midrule
Deterministic (baseline) & 36.5 & 100\% & Baseline & Plan exactly \\
$\pm 10\%$ stochastic & 37.8 & 98\% & +3.4\% & Add 5\% buffer \\
$\pm 20\%$ stochastic & 39.7 & 94\% & +8.8\% & Add 10\% buffer \\
$\pm 30\%$ stochastic & 42.8 & 87\% & +17.3\% & Add 15\% buffer \\
\bottomrule
\end{tabular}
\end{tabular}

\textbf{Key Insights:}
\begin{itemize}
\item Model handles moderate demand variability ($\pm 10\%$) well
\item High demand uncertainty requires 10-15\% safety stock buffer
\item Vehicle capacity constraints become binding with demand uncertainty
\end{itemize}

\subsection{Algorithm Sensitivity Analysis}

\subsubsection{Scenario Sampling Size}

We test impact of number of scenarios $S$ used in optimization:

\begin{table}[H]
\centering
\caption{Sensitivity to Number of Scenarios $S$}
\begin{tabular}{lcccc}
\toprule
Scenarios $S$ & Time (hrs) & Reliability Estimate & Runtime (sec) & Sampling Error \\
\midrule
20 & 34.2 & 0.91 & 8.7 & 4.1\% \\
50 & 35.7 & 0.90 & 23.4 & 2.2\% \\
100 (baseline) & 36.5 & 0.89 & 67.8 & 1.3\% \\
200 & 36.8 & 0.89 & 142.3 & 0.8\% \\
500 & 36.9 & 0.89 & 387.6 & 0.4\% \\
\bottomrule
\end{tabular}
\end{tabular}

\textbf{Key Insights:}
\begin{itemize}
\item $S = 100$ provides good balance (sampling error $< 2\%$)
\item Diminishing returns beyond $S = 200$ (computational cost outweighs accuracy)
\item Recommended: $S = 100$ for planning, $S_{\text{val}} = 1000$ for validation
\end{itemize}

\subsubsection{ALNS Operator Weights}

We test different initial operator weight configurations:

\begin{table}[H]
\centering
\caption{ALNS Operator Weight Sensitivity}
\begin{tabular}{lcccc}
\toprule
Weight Configuration & Final Solution & Iterations to Converge & Runtime & Recommendation \\
\midrule
Uniform (equal) & 37.2 & 8,450 & 67.8 & Baseline \\
Greedy-biased & 38.9 & 6,230 & 52.1 & Faster but worse \\
Explorative-biased & 36.8 & 12,780 & 98.3 & Better but slower \\
Adaptive (our method) & \textbf{36.5} & 8,450 & 67.8 & \textbf{Recommended} \\
\bottomrule
\end{tabular}
\end{tabular}

\textbf{Key Insights:}
\begin{itemize}
\item Adaptive mechanism automatically learns effective operators
\item Uniform initialization works well; adaptive mechanism optimizes search
\item Greedy bias risks local optima; explorative bias wastes computation
\end{itemize}

\subsection{Global Sensitivity Analysis (Sobol Indices)}

We perform variance-based global sensitivity analysis to identify most influential parameters:

\begin{table}[H]
\centering
\caption{Sobol Sensitivity Indices}
\begin{tabular}{lcc}
\toprule
Parameter & First-Order Index ($S_i$) & Total Effect Index ($S_{Ti}$) \\
\midrule
Reliability threshold $\alpha$ & 0.34 & 0.41 \\
Fleet size $K$ & 0.28 & 0.35 \\
Vehicle capacity $C$ & 0.18 & 0.22 \\
Demand variability & 0.12 & 0.16 \\
Travel time correlation $\rho$ & 0.08 & 0.11 \\
\bottomrule
\end{tabular}
\end{table}

\textbf{Interpretation:}
\begin{itemize}
\item Reliability threshold $\alpha$ and fleet size $K$ are most influential (34\% and 28\% of output variance)
\item Total effect indices confirm: no strong interaction effects ($S_{Ti} \approx S_i$)
\item Recommendation: Focus data collection on estimating these parameters accurately
\end{itemize}

\subsection{Stress Testing: Boundary Conditions}

We test model performance at boundary conditions:

\subsubsection{Minimum Reliability Requirement}

What if reliability requirement increases to $\alpha = 0.99$?
\begin{itemize}
\item Solution exists but requires 53\% more time
\item Cost increases by 47\%
\item Practical recommendation: $\alpha = 0.85$-$0.90$ optimal range
\end{itemize}

\subsubsection{Maximum Demand Surge}

What if demand suddenly increases by 50\% (emergency scenario)?
\begin{itemize}
\item Current fleet can handle 30\% surge
\item Beyond 30\%, need contingency plan (rent vehicles, request international aid)
\item Model identifies prioritization strategy (serve critical locations first)
\end{itemize}

\subsection{Robustness Summary}

Our sensitivity analysis demonstrates:
\begin{enumerate}
\item \textbf{Parameter robustness:} Model performs well across realistic parameter ranges
\item \textbf{Assumption robustness:} Solutions remain effective under moderate violations
\item \textbf{Algorithm robustness:} Adaptive mechanism automatically tunes search
\item \textbf{Identified vulnerabilities:} Strong correlation ($\rho > 0.6$) and sudden non-stationary changes require adaptive re-optimization
\end{enumerate}

These insights inform model evaluation in Section \ref{sec:evaluation} and practical implementation guidelines in Section \ref{sec:conclusion}.
